\subsection{Case Study}
\label{subsec:case-study}

To make the interaction between chain ambiguity and user intent drift concrete, Table~\ref{tab:case-study} traces a Joint anchor through its chain and drift phases. In the chain phase, the models must determine whether a \emph{disappointing} season means an overall losing record or poor away-game performance, and then decide whether the record should reflect the team's full season or the coach's stint. Both models identify the first layer but fail to surface the dependent record-scope ambiguity, applying \texttt{coaches.won/lost} instead of the gold \texttt{teams.won/lost}. In the drift phase, the user removes the 2000--2010 restriction. Both models correctly delete the year predicate (DRR=True), yet retain the unresolved record-scope error, yielding preEX=False and postEX=False. Because this constraint-removal drift leaves the record definition unchanged, the post-drift failure can be cleanly traced to the chain phase. This example shows that resolving $L1$ and correctly applying the drift are not sufficient for end-to-end correctness: the unresolved $L2$ ambiguity invalidates both queries.